% Notation macros. Fully portable: only amsmath/amssymb are assumed.
% Project-specific notation: see knowledge/writing/latex-conventions.md (section 4) and add it below.

% Number sets and expectation
\newcommand{\N}{\mathbb{N}}
\newcommand{\Z}{\mathbb{Z}}
\newcommand{\R}{\mathbb{R}}
\newcommand{\E}{\mathbb{E}}

% Operators
\newcommand{\OPT}{\mathrm{OPT}}
\newcommand{\diam}{\operatorname{diam}}
\newcommand{\poly}{\operatorname{poly}}

% Paired delimiters: \abs{x}, \ceil{x}, \floor{x}, \set{x}.
% With mathtools loaded (preamble/packages.tex) the starred forms \abs*{...} scale automatically;
% without it the plain definitions give the same output as the unstarred mathtools forms.
\makeatletter
\@ifpackageloaded{mathtools}{%
  \DeclarePairedDelimiter{\abs}{\lvert}{\rvert}%
  \DeclarePairedDelimiter{\ceil}{\lceil}{\rceil}%
  \DeclarePairedDelimiter{\floor}{\lfloor}{\rfloor}%
  \DeclarePairedDelimiter{\set}{\{}{\}}%
}{%
  \newcommand{\abs}[1]{\lvert #1\rvert}%
  \newcommand{\ceil}[1]{\lceil #1\rceil}%
  \newcommand{\floor}[1]{\lfloor #1\rfloor}%
  \newcommand{\set}[1]{\{#1\}}%
}
\makeatother

% Names of computational problems in small caps: \problemname{Min Cut-Path}
\newcommand{\problemname}[1]{\textsc{#1}}
